\section{Conclusion and Future Work}\label{sec:conclusion}

This thesis analysed a two-class non-preemptive priority $M/M/1$ queue under two waiting-room mechanisms ---jockeying and abandonment. It obtained the exact stationary bivariate PGF $P(x,y)$ in closed or integral form for five models: the baseline Model-$A$, one-way ($1\to2$) jockeying Model-$B_2$, class-$1$ abandonment Model-$C_2$, and the head-of-line variants Model-$\BH$ and Model-$\CH$.

One common property is shared by all solved models: the form of the fundamental equation that the PGF satisfies ---algebraic, or carrying derivative terms--- determines the tractability of the model with the methods employed in this work. As a result, the collection of closed-form solutions is smaller than the model hierarchy might suggest: derivative terms in $y$ complicate the problem, as one would need more advanced mathematical tools to determine the boundary function $P_y(y)$.


\subsection{The obstruction and the mechanism dichotomy}



Read as a competition between an algebraic coefficient and two derivative terms, the fundamental equation exposes a clean dichotomy. Where the equation is solvable, the unknown boundary function $P_y(y)$ is determined by exactly one of four mechanisms. The first is evaluation at a kernel root lying strictly inside the unit disc ---the purely algebraic route of the baseline Model-$A$. The second is interior analyticity: when one-way jockeying survives, the factor $(x-y)$ places the operative singularity at the diagonal point $x=y$ inside the disc, and the boundary function is fixed by requiring analyticity there (Model-$B_2$). The third is boundary boundedness: when class-$1$ abandonment survives, the factor $(x-1)$ moves the singularity to the unit-circle boundary $x=1$, where mere boundedness of the PGF suffices to determine the unknown (Model-$C_2$). The fourth is head-of-line algebraic collapse: restricting a mechanism to the head-of-line customer replaces the length-proportional rate with a constant one, eliminates the derivative term in the equation, and recovers the purely algebraic form of Model-$A$, solvable by the same kernel argument (Models~$\BH$ and~$\CH$).

The unsolvable cases ---Model-$B$ and Model-$X$--- make structural limits apparent, as the derivative terms in $y$ make the task of determining $P_y(y)$ significantly more challenging than it already was. Model-$B$ requires solving a Volterra-type integral equation whose forcing term depends on the boundary function itself; Model-$X$ further compounds this difficulty with non-linear characteristic curves, although we have not determined the exact difficulty or feasibility of the problem. The hierarchy is therefore not a loose collection of cases, but a complete enumeration: four mechanisms that enable exact solutions and two obstructions that prevent them.

The two mechanisms themselves act in opposite directions. Jockeying \emph{redistributes} delay at fixed total work: relocating a waiting customer conserves the total count, so $\mathbb{E}[N]$, $\pi_0$ and $\pi(0,0)$ are identical to the baseline for every $\gamma_1$. The mechanism's entire effect is to move delay between classes ---the priority premium $\mathbb{E}[W_2]/\mathbb{E}[W_1]$ inflating to $\premiumBtwoNinety$ at $\rho=0.90$ against the Model-$A$ value $\premiumANinety$.

Abandonment instead \emph{reduces} the load: the idle probability rises, every queue shortens, and total waiting times fall.

\subsection{Limitations}

Several assumptions are intrinsic to the $M/M/1$ framework itself. All results rest on exponential interarrival, service, jockeying, and abandonment times ---an idealisation that is most restrictive for the service times. Allowing class-dependent rates $\mu_1 \neq \mu_2$ would invalidate the analysis at its foundation, since the state-space construction and the definition of the limiting probabilities both rely on the server being indifferent to the class in service. Relaxing exponentiality entirely to general service times further compounds the difficulty; the intractability that already forces de~Waal's approximate $M/G/1$ treatment illustrates how quickly the problem ceases to admit closed-form solutions.

The partial-PGF result on the alternative state space $\widetilde{S}$ is an approximation, not a theorem, and should be read as a documented negative result. The rationale was as follows. Cohen~\cite{cohen1956delay} applied a partial generating function on the standard state space $(n_1,n_2)$ with respect to the class-2 variable, leaving the tractable class-1 subsystem untouched and transforming only in the harder direction. In the presence of jockeying, neither class has simple dynamics individually, but the total queue length $N = N_1 + N_2$ does. The alternative state space $(n_2, n)$ was therefore introduced to exploit this total-count structure. The attempt does not succeed: rather than simplifying the problem, the change of coordinates makes the balance equations harder even for Model-$A$.

\subsection{Future work}\label{sec:future_work}

Several extensions follow directly from the mathematical machinery used:
\begin{enumerate}
  \item \emph{Head-of-line versions of the remaining models.} The head-of-line models could be extended with bidirectional active mechanisms ---namely Model-$B^\mathrm{H}$ for jockeying and Model-$C^\mathrm{H}$ for abandonment--- rather than limiting them to strictly one direction, as the derivative terms in the fundamental equation would also disappear. Even a head-of-line model carrying both active mechanisms at the same time, Model-$X^\mathrm{H}$, would remain algebraic. These extensions appear to be within reach of the tools already developed here.
  \item \emph{The $k$-customer head-of-line ladder.} The head-of-line variants let only the single customer at the head of the queue for class~$1$ jockey or abandon; the natural extension lets the first $k$ waiting class-$1$ customers do so, jockeying at rate $\gamma_1\min(n_1,k)$ (and likewise abandoning at rate $\theta_1\min(n_1,k)$). Writing $\min(n_1,k)=\sum_{j=1}^{k}\mathbf 1_{\{n_1\ge j\}}$, the bounded rate keeps the fundamental equation \emph{algebraic} at the potential price of $k$ unknown boundary functions. This route has not been explored in this work, so it is unclear what it could yield. One would expect that such models would converge to Model-$B_2$ and Model-$C_2$ as $k\to\infty$.
  \item \emph{Numerical solution of the Model-$B$ Volterra equation.} Unlike Model-$X$, Model-$B$ yields a first-kind Volterra equation with an \emph{explicit} kernel and a \emph{known} forcing term. Whether its solution can be characterised further ---in closed form or through numerical approximation--- is left open; either route lies beyond the scope of this thesis.
  \item \emph{Multi-server extension.} Carrying the mechanisms to $M/M/c$ would connect this single-server line to the multi-server setting. With a common service rate for both classes, the non-preemptive priority $M/M/c$ queue admits an exact waiting-time analysis~\cite{kella1985waiting}, and exact methodology for two priority classes exists even with class-dependent rates under a preemptive discipline~\cite{wang2015mmc}, so exactness need not be surrendered at the outset; the open question is whether jockeying and abandonment can be added while preserving it. This would also connect to the (approximate) proactive-scheduling results of~\cite{huchandong2022proactive}.
  \item \emph{Derive mean performance metrics.} The closed-form PGFs obtained for Model-$B_2$, Model-$C_2$, Model-$\BH$, and Model-$\CH$ contain the full distributional information, but extracting mean queue lengths and waiting times requires differentiating at $(1,1)$, which in Model-$B_2$ and Model-$C_2$ involves boundary functions given as integrals. A systematic differentiation ---applying L'H\^opital where necessary at the singularity $x=1$--- would yield closed-form expressions for $\mathbb{E}[N_1]$, $\mathbb{E}[N_2]$, and, via Little's law, $\mathbb{E}[W_1]$ and $\mathbb{E}[W_2]$, replacing the numerical estimates reported in the results chapter with exact ones.
  \item \emph{Calibration.} Fitting $\gamma_1$, $\theta_1$ and the loads to real elderly-care waiting-list data~\cite{zenios1999modeling} would test whether the regimes identified here accurately describe operational reality. Since the transplant waiting-list model of~\cite{zenios1999modeling} differs from ours in several respects, such a calibration would likely require further model extensions rather than a direct fit of $\gamma_1$ and $\theta_1$.
\end{enumerate}
